% Supplemental proof supplied to the C owner.  This file has no automatic
% input in AB_push.tex; its C-number is to be assigned after independent review.
\begin{lemma}[AB--C comparison: the mixed class determines the square]
Let $\pi:V\to C$ be a compact oriented smooth genus-one Lefschetz
fibration over a closed oriented surface, all of whose singular fibres
are reduced irreducible positive nodal fibres.  Suppose its smooth
fibres have a smoothly varying elliptic complex structure and affine
holomorphic transition maps.  Let $J\to C$ be its marked Jacobian
fibration.  Let $L$ be a complex line bundle on $V$, of degree zero on
every fibre, whose fibrewise holomorphic structures define a specified
smooth section $P$ of the relative degree-zero Picard family on the
regular locus.  Suppose this section extends as the same section for
$J$, and let $L_J$ be a line bundle on $J$ representing it, also of
fibre degree zero.  Then
\[
 \langle c_1(L)^2,[V]\rangle
       =\langle c_1(L_J)^2,[J]\rangle.
\]
It suffices that the assertion about the fibrewise Picard section be
checked off disjoint nodal discs.  No projectivity, multisection,
holomorphic total-space torsor, or holomorphic dependence of an affine
origin on the base is required.
\end{lemma}
\begin{proof}
Put $j:C^\circ\hookrightarrow C$ for the complement of the critical
values and $\Lambda=j_*R^1\pi^\circ_*\Z$.  The local reduced $I_1$
model has
\[
 R^0\pi_*\Z=\Z,\qquad R^1\pi_*\Z=\Lambda,
 \qquad R^2\pi_*\Z=\Z,
 \qquad R^q\pi_*\Z=0\quad(q>2).
\]
The constant generator in degree two is normalized by the oriented
fibre.  The same assertions and the same cup-product pairing
$\Lambda\otimes\Lambda\to\Z$ hold for $J$.  Indeed a nodal fibre
is a torus with one primitive cycle collapsed.  Its degree-one
cohomology is the invariant subgroup of the regular local system,
and its fundamental class is the degree-one specialization of the
nearby fundamental class.  This gives the displayed sheaves directly,
including their integral normalization.

The Leray filtration has
\[
 F^1H^2(V;\Z)/F^2H^2(V;\Z)
       =H^1(C,\Lambda),\qquad
 H^4(V;\Z)=H^2(C;\Z).
\]
There is no differential into or out of $E_2^{1,1}$, since the base
has cohomological dimension two, and the only term of total degree
four is $E_2^{2,2}$.  Fibre degree zero puts $e=c_1(L)$ in $F^1$.
Multiplicativity now says that its square is the product of its
class in $H^1(C,\Lambda)$ with itself, using the fibre intersection
pairing.  Changing $e$ by an element of $F^2$ cannot change that
square: its product with $F^1$ lies in $F^3=0$.

We spell out why the mixed class is the one specified by $P$.
Over a surface with boundary obtained by deleting the nodal discs,
the affine elliptic torsor has a smooth section.  The obstruction is
in $H^2$ of its lattice local system; this group is zero because
that base retracts onto a graph.  Equivalently, the exact sequence
of sheaves of smooth sections of the Lie bundle and the torus bundle
has fine middle sheaf and gives precisely this obstruction group.
Translation by such a section identifies the torsor with its
Jacobian.  On a sufficiently small base chart, normalize $L$ at the
chosen origin.  Its fibrewise holomorphic isomorphism class is then
the Poincar\'e line specified by the smooth section $P$.
Two such normalized identifications differ by a function of the
base, because an automorphism of a holomorphic line bundle on a
compact connected elliptic curve is constant.  Thus their possible
remaining discrepancy is a line from the base.  On this graph-like
base its first Chern class is zero.  In sheaf language, the mixed
class is the connecting class of $P$ for the relative exponential
sequence; this is the same class for the torsor and its Jacobian.

This equality on $C^\circ$ determines the global mixed class:
the degree-one edge sequence for $j$ gives an injection
\[
 H^1(C,j_*\Lambda^\circ)
       \longrightarrow H^1(C^\circ,\Lambda^\circ).
\]
Therefore $e$ and $c_1(L_J)$ have equal mixed classes in the
canonically identified $H^1(C,\Lambda)$.  Their squares agree by
the preceding Leray product calculation.  This also explains why
the possible absence of a global affine origin does not affect the
answer.
\end{proof}

\begin{lemma}[AB--C comparison: nodal excision on a page]
Let $M\to\Sigma$ be the corresponding compact genus-one Lefschetz
surface with nonempty base boundary, and remove the interiors of
the inverse images $E_i$ of disjoint small discs around its nodal
values.  Then
\[
 H^2(M;\Z)\longrightarrow
 H^2\bigl(M\setminus\textstyle\bigcup_i\operatorname{int}E_i;\Z\bigr)
\]
is injective.  On a closed base, its kernel instead consists of
multiples of the fibre class.
\end{lemma}
\begin{proof}
Excision and Poincar\'e--Lefschetz duality identify each local
degree-two group with
\[
 H^2(E_i,\partial E_i;\Z)\simeq H_2(E_i;\Z)=\Z.
\]
The pullback of the relative base-disc generator is primitive:
its Poincar\'e dual is the regular fibre, and that fibre specializes
with multiplicity one to the reduced irreducible nodal fibre,
which generates $H_2(E_i;\Z)$.  The image in global cohomology is
therefore the pullback of the point class of the base.  This is
zero if $\partial\Sigma\ne\varnothing$; on a closed base all such
images are multiples of the same fibre class.  The cohomology
sequence of the pair proves the two assertions.  In particular,
an additional absolute mixed Chern class cannot be introduced by
changing a nodal clutch while leaving the exterior bundle fixed.
\end{proof}

\begin{proposition}[AB--C comparison: the actual line has square $-2$]
Assume the actual marked constructions of C.3.13, C.3.17 and
C.3.18--C.3.20.  With the Lefschetz orientation used there, the
actual circle bundle on the actual bounding four-manifold has
\[
 s_\delta=-2.
\]
This statement uses the actual semiabelian transition maps and is
independent of the auxiliary genus-three plumbing of B.3.7.
\end{proposition}
\begin{proof}
We first identify the base and its Jacobian.  For
$F_\epsilon=(1-\epsilon)z_2^3-z_1^4$, where $0<\epsilon<1$,
the positive weighted normalization of a page gives
\[
 a=\frac{z_1}{F_\epsilon^{1/4}},\qquad
 b=\frac{(1-\epsilon)^{1/3}z_2}{F_\epsilon^{1/3}},
 \qquad b^3=a^4+1.
\]
The compactification is the genus-three curve $C$ with one point
$\infty_C$, where $a$ and $b$ have pole orders $3$ and $4$.
The actual elliptic-base map is
\[
 t=\frac{a^4+1}{1+\epsilon a^4}.
\]
Its ramification orders are $3,4,12$ over $0,1,1/\epsilon$;
the twelve preimages of the original cusp $t=\infty$ are unramified.
This is also the degree-twelve cover
\[
 y^{12}=\frac{t^4}{(t-1)^3(t-1/\epsilon)}
\]
after a constant rescaling of $y$.  The Riemann--Hurwitz sum is
$-24+8+9+11=4$.

Write $\lambda=1-\epsilon$ and $h=1+\epsilon a^4$.  On the
pulled-back Weierstrass equation the change
\[
 x=(ab^2/h)^2X,\qquad y=(ab^2/h)^3Y
\]
gives its minimal Jacobian model
\[
 J:\quad Y^2=X^3-3\lambda bX+2\lambda^2a^2,
 \qquad P=(0,-\sqrt2\lambda a).
\]
Its discriminant is $1728\lambda^3h$.  The four simple roots
of $h$ each have three nonzero $b$-values.  There are exactly
twelve $I_1$ fibres.  At $\infty_C$ the Weierstrass coefficients
have respective pole orders $4$ and $6$, and the discriminant
has pole order $12$ with nonzero leading coefficient.  Thus
the fundamental line is $\mathcal O_C(\infty_C)$, the fibre
there is smooth, and $\chi(\mathcal O_J)=1$.
The section $P$ is affine at every finite point.  At infinity,
after the minimal scaling, its $X$ coordinate is zero and its
$Y$ coordinate has a finite nonzero limit.  It is therefore
disjoint from $O$.  There are no reducible-fibre corrections.
The section normal bundle has degree $-1$, so
\[
 c_1\bigl(\mathcal O_J(P-O)\bigr)^2
       =P^2+O^2-2P\cdot O=-1-1=-2.
\]
Equivalently this is minus the height $2$ of the pulled-back
section; compare \citet[Theorem 8.6]{Shioda90}.

We now check the actual bundle, rather than inferring it from
the Jacobian.  On a selected ordinary page chart let $g$ be
the real $\gamma$ coefficient.  In the marked period coordinates
put
\[
 \xi=\zeta_1-6\mu g=\tau u+w,
 \qquad \eta=\zeta_2-\beta g=\mu u+d.
\]
Then $W=\exp(2\pi i\eta)$ has exactly the factors
\[
 (\xi,W)\sim(\xi+1,W),\qquad
 (\xi,W)\sim(\xi+\tau,e^{2\pi i\mu}W).
\]
The meromorphic theta quotient
$\vartheta(\xi-\mu)/\vartheta(\xi)$ has these factors and divisor
$P-O$.  Hence this is the actual Poincar\'e line of $P$ on each
chart, with its fibrewise holomorphic structure and its actual
smooth unit-circle bundle.  The two ordinary marking changes are
\[
 R_1=\begin{pmatrix}-1/\tau&0\\(1-\mu)/\tau&1\end{pmatrix},
 \qquad
 R_2=\begin{pmatrix}1/\tau&0\\-\mu/\tau&1\end{pmatrix}.
\]
They satisfy $R_j\Pi=\Pi' A_j$, preserve the $\delta$ column,
and preserve the $\gamma$ character.  They are triangular;
after subtracting $g$ times the first period their changes in
$\xi$ are affine holomorphic in the elliptic fibre coordinate,
and their changes in $W$ are the corresponding holomorphic
line lifts.  The additional translations can depend smoothly,
and need not depend holomorphically, on the base.  This does
not change the fibrewise Picard family.  Explicitly,
\[
 \mu'=-\mu/\tau+(1-\tau')\quad(j=1),
 \qquad \mu'=\mu/\tau+1\quad(j=2),
\]
so the resulting relative Picard section is the globally marked
$P$.  Integral changes of the lift of $g$ give the specified
semiabelian lift of translation by $6P$; translation preserves
degree-zero Picard classes.  On a finite root chart the additional
affine translation has the same property.  The page cover unwraps
the orders $3$ and $4$, so the two finite root charts extend as
smooth affine elliptic bundles there.  This checks equality as
a smooth section of the relative Picard family, including its
mixed class, rather than just fibrewise topological degree zero.

Remove the twelve entire nodal inverse-image discs from the
actual page surface.  On the remainder all these identifications
are the original marked coordinates.  The cusp transport of
C.3.17 keeps every real period coefficient fixed on its outer
collar, so it retains the same local systems and line comparison
there.  The preceding nodal-excision lemma shows that no
unseen absolute Chern class is introduced inside the twelve
nodal pieces.  Alternatively, the injectivity for $j_*$ in the
first lemma shows that their global mixed class is determined
by this exterior computation.

Finally cap the regular page boundary by a marked
$T^2\times D^2$.  To make this literal, first remove a smaller
product collar of the page boundary.  Its interior is the
punctured algebraic curve above; the remaining end is a disc
about $\infty_C$ on which the Jacobian is smooth.  The smooth
affine torus bundle over the annulus has a section, since its
twisted $H^2$ vanishes.  The actual Picard section is the
restriction of $P$ there, and hence extends over the disc.
Choose the line on this regular cap representing that extension.
Its restriction can be glued by a fibrewise holomorphic line
isomorphism; all further factors are base factors.  Such a
degree-zero line over $T^2\times D^2$ is topologically trivial.
Use one of its smooth trivializations on the common boundary.
The relative square on the original page equals the absolute
square on the capped Lefschetz surface.  Changing from this
boundary trivialization to the actual marked one cannot change
the square, by the radical argument proved in C.3.20.

The first lemma applies to this capped actual affine Lefschetz
surface and to $J$.  It identifies their squares and gives
$s_\delta=-2$ with the stated orientation.
\end{proof}

\begin{remark}[Audit and finite certificates]
The assertions about the actual page, the marked cusp transport
and its disjointness from the fixed sphere are the previously
proved C constructions, not new consequences of a height formula.
The exact script \texttt{ab\_push\_actual\_line\_comparison.py}
checks both period covariance identities, their fixed character
and circle column, the pulled-back minimal equation and section,
and the discriminant.  Its output is saved beside the script.
The proof deliberately makes no assertion that the affine
coordinates $\xi$ depend holomorphically on the page parameter.
\end{remark}
